% Subsection 3.2.2: 解答与测试一致性校验

格式层只能保证字段形态正确，无法保证 \texttt{solution} 与 \texttt{test} 语义匹配；KodCode 中确实存在解答与测试版本不同步、依赖未声明、或测试覆盖的函数名与 \texttt{test\_info} 不一致等问题。为此 \texttt{validate\_code\_with\_test.py} 实现了轻量执行器 \texttt{test\_solution\_code}，在数据清洗与 pass@k 评测两个场景复用同一套判定逻辑，避免口径漂移。

\paragraph{（1）执行策略}
校验以「子进程 + 标准库」为原则，不引入第三方 runner。脚本将 \texttt{solution} 与 \texttt{test} 写入临时目录并在独立子进程中运行：进程隔离保证被测代码的全局副作用（如修改 \texttt{sys.path}、覆盖内置名）不污染主进程，超时上限用于杀掉死循环与指数回溯等病态用例。

\paragraph{（2）判定语义}
返回 \texttt{(ok, err)} 二元组：\texttt{ok=True} 当且仅当全部断言通过且进程退出码为 $0$；测试失败、导入错误、\texttt{SyntaxError}、超时等一切情形均记为失败并回填简短错误消息。这一「唯一正确定义」贯穿数据清洗与评测两个阶段，避免口径漂移。

\paragraph{（3）一致性收益}
最终保留的样本满足「官方 \texttt{solution} 能通过其自带 \texttt{test}」这一硬约束，为后续训练与评测提供严格上界。评测集上参考解答 pass@1 实测为 $0.996$（见 3.4 节），与 $1$ 的微小差距几乎全部来自执行非确定性因素（如浮点精度、系统时钟）而非数据质量缺陷。